\documentclass{amsart}
\usepackage{amsmath,amssymb,amsthm,hyperref,stmaryrd}
\usepackage{algorithm}
\usepackage{algpseudocode}

\theoremstyle{plain}
\newtheorem{theorem}{Theorem}
\newtheorem{lemma}[theorem]{Lemma}
\newtheorem{proposition}[theorem]{Proposition}
\newtheorem{corollary}[theorem]{Corollary}
\theoremstyle{definition}
\newtheorem{definition}[theorem]{Definition}
\newtheorem{remark}[theorem]{Remark}

\newcommand{\To}{\rightarrow}
\newcommand{\co}{\prec}            % co-implication / subtraction
\newcommand{\bigC}{\mathbf{2C}}
\newcommand{\Ci}{\mathbf{2C}^{-{\sim}}}        % system (i): no strong negation
\newcommand{\Cii}{\mathbf{2C}^{-{\sim},-c}}    % system (ii): no strong negation, no constants
\newcommand{\fp}{\Vdash^{+}}
\newcommand{\fm}{\Vdash^{-}}
\newcommand{\proves}{\vdash}
\newcommand{\refutes}{\dashv}

\begin{document}

\title{Bilateral fragments of the bi-connexive logic $\mathbf{2C}$:\\
deleting the strong negation and the propositional constants}
\author{solver-agent}
\date{}
\maketitle

\begin{abstract}
We analyse the two bilateral propositional systems obtained from the
bi-connexive logic $\bigC$ by (i) deleting the strong (toggle) negation while
keeping $\top,\bot$, and (ii) additionally deleting $\top,\bot$. For each system
we give a sound and complete bilateral natural--deduction calculus, a relational
and an algebraic (twist--structure) semantics, and we settle the status of the
connexive theses, of negation--inconsistency, of non--triviality, and the
definability/undefinability of the absent operators. The decisive structural
fact is a \emph{separation invariant}: in the negation--free language the
support--of--truth of a formula depends only on the positive valuation of the
atoms. From it we derive (a)~the strong negation is undefinable in both systems;
(b)~the connexive theses, which are statements \emph{about} the strong negation,
collapse; (c)~system~(i) is negation--inconsistent but non--trivial, witnessed by
$\bot\To p$; and (d)~system~(ii) is negation--consistent, because no
constant--free formula is ever refuted.
\end{abstract}

\section{The base logic $\bigC$ and its bilateral semantics}\label{sec:base}

We work with Wansing's bilateral formulation. A \emph{bi--consequence} system has
two relations: a \emph{provability} relation $\proves$ and a \emph{refutability}
relation $\refutes$. Semantically these are read off a single Kripke model
through a pair of support relations $\fp$ (``is verified at'') and $\fm$ (``is
falsified at'').

\begin{definition}[Models]\label{def:model}
A \emph{model} is a triple $\mathfrak{M}=(W,\le,V)$ where $(W,\le)$ is a
non--empty poset (a reflexive--transitive frame) and
$V=(V^{+},V^{-})$ assigns to each atom $p$ two \emph{$\le$--upward closed}
sets $V^{+}(p),V^{-}(p)\subseteq W$ (\emph{heredity}: if $w\le v$ and $w\in
V^{\pm}(p)$ then $v\in V^{\pm}(p)$). No relation between $V^{+}(p)$ and
$V^{-}(p)$ is imposed.
\end{definition}

\begin{definition}[Support, full language]\label{def:support}
The full language of $\bigC$ is built from atoms with
$\wedge,\vee,\To$, the strong negation ${\sim}$, and the constants
$\top,\bot$. The two support relations are fixed by mutual recursion:
\begin{align*}
w\fp p &\iff w\in V^{+}(p), & w\fm p &\iff w\in V^{-}(p),\\
w\fp \top &\iff \text{always}, & w\fm \top &\iff \text{never},\\
w\fp \bot &\iff \text{never}, & w\fm \bot &\iff \text{always},\\
w\fp A\wedge B &\iff w\fp A \text{ and } w\fp B,
 & w\fm A\wedge B &\iff w\fm A \text{ or } w\fm B,\\
w\fp A\vee B &\iff w\fp A \text{ or } w\fp B,
 & w\fm A\vee B &\iff w\fm A \text{ and } w\fm B,\\
w\fp A\To B &\iff \forall v\ge w\,(v\fp A\Rightarrow v\fp B),
 & w\fm A\To B &\iff \forall v\ge w\,(v\fp A\Rightarrow v\fm B),\\
w\fp {\sim}A &\iff w\fm A, & w\fm {\sim}A &\iff w\fp A .
\end{align*}
We write $\mathfrak{M}\models^{+}A$ if $w\fp A$ for all $w\in W$, and
$\mathfrak{M}\models^{-}A$ if $w\fm A$ for all $w\in W$; and $\models^{+}A$
(resp.\ $\models^{-}A$) if this holds in every model. The two clauses for
$\To$ that mention the falsity component are the \emph{connexive twist}: the
falsity of $A\To B$ propagates the falsity of the \emph{consequent} under the
truth of the antecedent. The strong negation ${\sim}$ is the toggle linking the
two fragments.
\end{definition}

\begin{lemma}[Heredity]\label{lem:hered}
For every formula $A$ of the full language and every model, the sets
$\{w: w\fp A\}$ and $\{w: w\fm A\}$ are $\le$--upward closed.
\end{lemma}
\begin{proof}
Induction on $A$. Atoms hold by Definition~\ref{def:model}; $\top,\bot$ are
constant hence upward closed. For $\wedge,\vee$ closure is preserved by finite
intersection/union of upward closed sets. For $\To$ both clauses are universal
quantifications over $\{v:v\ge w\}$ of an upward persistent matrix, hence are
themselves upward closed because $\ge$ is transitive. For ${\sim}$ the two
components are swapped, so closure of one gives closure of the other.
\end{proof}

\begin{proposition}[Connexive theses and bi--connexivity of $\bigC$]\label{prop:base}
In every model,
\[
\models^{+}{\sim}(A\To{\sim}A),\quad
\models^{+}{\sim}({\sim}A\To A),
\]
\[
\models^{+}(A\To B)\To{\sim}(A\To{\sim}B),\quad
\models^{+}(A\To{\sim}B)\To{\sim}(A\To B),
\]
and dually $\models^{-}(A\To{\sim}A)$, $\models^{-}({\sim}A\To A)$;
i.e.\ Aristotle's and both directions of Boethius' theses hold in the
provability fragment, and their duals in the refutability fragment.
\end{proposition}
\begin{proof}
By Definition~\ref{def:support}, $w\fp{\sim}(A\To{\sim}A)$ iff $w\fm (A\To{\sim}A)$
iff $\forall v\ge w\,(v\fp A\Rightarrow v\fm{\sim}A)$ iff $\forall v\ge
w\,(v\fp A\Rightarrow v\fp A)$, which is a tautology. The case of
${\sim}({\sim}A\To A)$ is symmetric, using $w\fm{\sim}A\iff w\fp A$.
For Boethius, $w\fp(A\To B)\To{\sim}(A\To{\sim}B)$ holds iff for all $v\ge w$
with $v\fp A\To B$ we have $v\fp{\sim}(A\To{\sim}B)$, i.e.\ $v\fm(A\To{\sim}B)$,
i.e.\ $\forall u\ge v\,(u\fp A\Rightarrow u\fm{\sim}B)$, i.e.\
$\forall u\ge v\,(u\fp A\Rightarrow u\fp B)$; and the last statement is exactly
$v\fp A\To B$, which holds by assumption. The other Boethius direction is the
same computation with the roles of $B$ and ${\sim}B$ exchanged. The refutability
duals follow because $w\fm(A\To{\sim}A)\iff\forall v\ge w\,(v\fp A\Rightarrow
v\fm{\sim}A)\iff\forall v\ge w\,(v\fp A\Rightarrow v\fp A)$, again a tautology.
\end{proof}

(These claims were also verified by exhaustive model checking over all reflexive
transitive frames with $\le 3$ worlds and all hereditary valuations.)

\section{The two reduced systems}\label{sec:systems}

\begin{definition}\label{def:reduced}
\emph{System $\Ci$} is the bi--consequence system in the language
$\{\wedge,\vee,\To,\top,\bot\}$ (strong negation deleted), with $\fp,\fm$ as in
Definition~\ref{def:support} restricted to these connectives. \emph{System
$\Cii$} is the further restriction to the constant--free language
$\{\wedge,\vee,\To\}$.
\end{definition}

Both are interpreted in exactly the models of Definition~\ref{def:model}; only
the stock of connectives changes. We write $\Vdash_{\Ci}^{+}A$, etc.

\subsection{The separation invariant}

\begin{lemma}[Separation invariant]\label{lem:sep}
Let $A$ be a formula of the language $\{\wedge,\vee,\To,\top,\bot\}$ (i.e.\
strong--negation--free). Then in every model the truth set $\{w:w\fp A\}$ is a
function of $\le$ and of the positive valuation $V^{+}$ alone; it does not depend
on $V^{-}$.
\end{lemma}
\begin{proof}
Structural induction on $A$. For atoms $w\fp p\iff w\in V^{+}(p)$, independent of
$V^{-}$. For $\top,\bot$ the truth set is $W$ resp.\ $\emptyset$, independent of
$V$. For $\wedge,\vee$ the truth set is the intersection resp.\ union of the
truth sets of the immediate subformulas, which by the induction hypothesis depend
only on $V^{+}$. For $\To$, the clause
$w\fp A\To B\iff\forall v\ge w\,(v\fp A\Rightarrow v\fp B)$ refers only to the
truth sets of $A$ and $B$, which by hypothesis depend only on $V^{+}$. As there
is no clause for ${\sim}$ in this language, the induction is complete.
\end{proof}

\begin{remark}\label{rem:nofalse}
The dual statement is \emph{false}: $w\fm A\To B$ uses $v\fp A$, so the falsity
set of a strong--negation--free formula generally depends on \emph{both}
components. This asymmetry is the engine of all that follows.
\end{remark}

\subsection{Undefinability of the strong negation}

\begin{theorem}[Undefinability of ${\sim}$]\label{thm:undefneg}
There is no formula $N(p)$ of the language of $\Ci$ (a fortiori none of $\Cii$)
such that for every model and world, $w\fp N(p)\iff w\fm p$ and $w\fm N(p)\iff
w\fp p$. Hence the strong negation is not definable in either reduced system,
even up to provable/refutable equivalence.
\end{theorem}
\begin{proof}
Suppose such an $N(p)$ existed. Consider two models on the single reflexive world
$W=\{w\}$ ($\le$ the identity), differing only in $V^{-}(p)$:
$\mathfrak{M}_{0}$ with $V^{+}(p)=\{w\}$, $V^{-}(p)=\emptyset$, and
$\mathfrak{M}_{1}$ with $V^{+}(p)=\{w\}$, $V^{-}(p)=\{w\}$. By
Lemma~\ref{lem:sep} the truth set of $N(p)$ is the same in $\mathfrak{M}_{0}$ and
$\mathfrak{M}_{1}$, since the two models share the same $V^{+}$ and the same
frame. In particular $w\fp N(p)$ has the same truth value in both. But the
defining property would force $w\fp N(p)\iff w\fm p$, which is \emph{false} in
$\mathfrak{M}_{0}$ ($w\not\fm p$) and \emph{true} in $\mathfrak{M}_{1}$
($w\fm p$). This contradiction shows no such $N$ exists.
\end{proof}

\subsection{Undefinability of the constants in $\Cii$}

\begin{theorem}[Undefinability of $\top,\bot$]\label{thm:undefconst}
Neither $\top$ nor $\bot$ is definable in $\Cii$: there is no constant--free
formula $T$ with $(w\fp T,w\fm T)=(\text{true},\text{false})$ in all models, and
none with $(\text{false},\text{true})$ in all models. Indeed the same holds even
if the strong negation were re--added but the constants kept out.
\end{theorem}
\begin{proof}
Take the single reflexive world $W=\{w\}$ with $V^{+}(p)=V^{-}(p)=\{w\}$ for
every atom $p$, so each atom satisfies $(w\fp p,w\fm p)=(\text{true},\text{true})$.
We claim every constant--free formula $A$ (allowing also ${\sim}$) satisfies
$(w\fp A,w\fm A)=(\text{true},\text{true})$. Induction: atoms hold by choice; for
$A\wedge B$, $w\fp A\wedge B$ is $\text{true}\wedge\text{true}$ and $w\fm A\wedge
B$ is $\text{true}\vee\text{true}$, both true; dually for $\vee$; for $A\To B$ on
the reflexive single world, $w\fp A\To B\iff(w\fp A\Rightarrow w\fp B)=
(\text{true}\Rightarrow\text{true})=\text{true}$ and $w\fm A\To B\iff(w\fp
A\Rightarrow w\fm B)=\text{true}$; for ${\sim}A$ the pair is swapped, staying
$(\text{true},\text{true})$. Thus no constant--free formula attains the pair
$(\text{true},\text{false})$ of $\top$ or $(\text{false},\text{true})$ of $\bot$
in this model, so neither constant is definable.
\end{proof}

\section{Proof theory: sound and complete bilateral calculi}\label{sec:proof}

We use a bilateral natural deduction in which every derivable judgement is a
\emph{signed} statement $A^{+}$ (``$A$ is proved'') or $A^{-}$ (``$A$ is
refuted''); assumptions are signed likewise, and $\proves,\refutes$ are the
single--conclusion consequence relations
$\Gamma\proves A:\!\!\iff \Gamma\Rightarrow A^{+}$ and
$\Gamma\refutes A:\!\!\iff\Gamma\Rightarrow A^{-}$.

\begin{definition}[Calculus $\mathsf{N}(\Ci)$]\label{def:Ncalc}
Sequents are $\Gamma\Rightarrow C^{s}$ with $s\in\{+,-\}$, $\Gamma$ a finite set
of signed formulas. The rules are:

\smallskip\noindent\emph{Structural:} identity $A^{s}\Rightarrow A^{s}$;
weakening; and cut on a signed formula
$\dfrac{\Gamma\Rightarrow A^{s}\quad \Delta,A^{s}\Rightarrow C^{t}}
       {\Gamma,\Delta\Rightarrow C^{t}}$.

\smallskip\noindent\emph{Constants:} $\Rightarrow\top^{+}$, $\Rightarrow\bot^{-}$,
and the explosion rules $\bot^{+}\Rightarrow C^{+}$, $\bot^{+}\Rightarrow C^{-}$,
$\top^{-}\Rightarrow C^{+}$, $\top^{-}\Rightarrow C^{-}$.

\smallskip\noindent\emph{Conjunction:}
$A^{+},B^{+}\Rightarrow(A\wedge B)^{+}$; $(A\wedge B)^{+}\Rightarrow A^{+}$,
$(A\wedge B)^{+}\Rightarrow B^{+}$;
$A^{-}\Rightarrow(A\wedge B)^{-}$, $B^{-}\Rightarrow(A\wedge B)^{-}$;
$\dfrac{\Gamma,A^{-}\Rightarrow C^{s}\quad\Gamma,B^{-}\Rightarrow C^{s}}
       {\Gamma,(A\wedge B)^{-}\Rightarrow C^{s}}$.

\smallskip\noindent\emph{Disjunction:} dual to conjunction (swap $+/-$ and
$\wedge/\vee$).

\smallskip\noindent\emph{Implication, positive part} (intuitionistic):
$\dfrac{\Gamma,A^{+}\Rightarrow B^{+}}{\Gamma\Rightarrow(A\To B)^{+}}$ and
$\dfrac{\Gamma\Rightarrow(A\To B)^{+}\quad\Delta\Rightarrow A^{+}}
       {\Gamma,\Delta\Rightarrow B^{+}}$, with the proviso (for the
$\To^{+}$--introduction and the $\To^{-}$--introduction below) that $\Gamma$
contains only positively signed formulas, encoding the heredity restriction.

\smallskip\noindent\emph{Implication, connexive negative part:}
$\dfrac{\Gamma,A^{+}\Rightarrow B^{-}}{\Gamma\Rightarrow(A\To B)^{-}}$ and
$\dfrac{\Gamma\Rightarrow(A\To B)^{-}\quad\Delta\Rightarrow A^{+}}
       {\Gamma,\Delta\Rightarrow B^{-}}$.
\end{definition}

\begin{definition}[Calculus $\mathsf{N}(\Cii)$]
Obtained from $\mathsf{N}(\Ci)$ by deleting the four constant rules and the two
introduction/elimination axioms for $\top,\bot$.
\end{definition}

\begin{theorem}[Soundness and completeness]\label{thm:sc}
For all finite $\Gamma$ and formulas $C$ in the respective language and
$s\in\{+,-\}$:
\[
\mathsf{N}(\Ci)\vdash \Gamma\Rightarrow C^{s}
\iff \Gamma\models C^{s},
\qquad
\mathsf{N}(\Cii)\vdash \Gamma\Rightarrow C^{s}
\iff \Gamma\models C^{s},
\]
where $\Gamma\models C^{+}$ means: in every model and world $w$, if $w\fp A$ for
all $A^{+}\in\Gamma$ and $w\fm A$ for all $A^{-}\in\Gamma$, then $w\fp C$; and
$\Gamma\models C^{-}$ is the analogous statement with $w\fm C$ in the conclusion.
\end{theorem}
\begin{proof}
\emph{Soundness.} By induction on derivations one checks every rule preserves the
semantic consequence $\models$, fixing a model and a world $w$. The constant and
Boolean rules are immediate from Definition~\ref{def:support}. The positive
implication rules are the standard intuitionistic ones; the restriction to
positive contexts in $\To$--introduction is exactly heredity
(Lemma~\ref{lem:hered}), which licenses passing to an arbitrary $v\ge w$. The
connexive negative rules are sound because the falsity clause for $\To$ has the
same universal--implicational shape as the truth clause, with $B^{-}$ in place of
$B^{+}$: from $\Gamma,A^{+}\Rightarrow B^{-}$ one gets, for any $v\ge w$
verifying the positive context and any $u\ge v$ with $u\fp A$, that $u\fm B$,
which is $v\fm A\To B$. Explosion for $\bot^{+}$ and $\top^{-}$ is sound because
no world verifies $\bot$ nor falsifies $\top$, so the premise set is
unsatisfiable.

\emph{Completeness.} Use the bilateral canonical model. Call a set $\Phi$ of
signed formulas \emph{prime--saturated} if it is deductively closed under
$\mathsf{N}$, does not derive every signed formula, satisfies the disjunction
property on the $+$--component and the dual conjunction property on the
$-$--component, and is consistent in the sense that it does not contain a signed
formula together with a derivation of the empty conclusion. By a standard
Lindenbaum construction every underivable sequent $\Gamma\Rightarrow C^{s}$
extends $\Gamma$ to a prime--saturated $\Phi$ omitting $C^{s}$. Let $W$ be the set
of prime--saturated sets ordered by inclusion of their $+$--parts together with
\emph{reverse} inclusion of their $-$--parts in the standard 2Int manner; set
$\Phi\fp p\iff p^{+}\in\Phi$, $\Phi\fm p\iff p^{-}\in\Phi$. Heredity holds by
construction. A truth lemma, proved by induction on formulas using the
saturation properties exactly as for intuitionistic logic on the positive side
and, on the negative side, using the connexive negative rules to handle $A\To B$
($\Phi\fm A\To B\iff(A\To B)^{-}\in\Phi$ is obtained from the negative
introduction/elimination pair), gives $\Phi\fp A\iff A^{+}\in\Phi$ and
$\Phi\fm A\iff A^{-}\in\Phi$. Then the canonical $\Phi$ omitting $C^{s}$
refutes the sequent, proving completeness. Deleting the constant rules and
constant vocabulary yields $\mathsf{N}(\Cii)$ over the constant--free models,
and the identical argument applies, the constant clauses now being vacuous.
\end{proof}

\begin{remark}[The positive fragment is intuitionistic]\label{rem:int}
By Lemma~\ref{lem:sep} the relation $\Gamma^{+}\models C^{+}$ (positive
assumptions only) coincides with intuitionistic consequence: the positive
clauses are the intuitionistic Kripke clauses and the negative valuation is
irrelevant. Hence the provability fragment of $\Ci$ is intuitionistic logic
$\mathbf{IPC}$, and that of $\Cii$ is its $\{\wedge,\vee,\To\}$--fragment
(positive intuitionistic logic). This was confirmed numerically: Peirce's law and
excluded middle fail in the $+$--reading, while $A\To(B\To A)$ and $A\To A$ hold.
\end{remark}

\section{Algebraic and relational semantics}\label{sec:alg}

\begin{definition}[Twist algebras]\label{def:twist}
Let $\mathbf H=(H,\wedge,\vee,\To_{H},0,1)$ be a Heyting algebra. The
\emph{connexive twist algebra} $\mathbf H^{\bowtie}$ has carrier $H\times H$,
where the first coordinate tracks truth and the second falsity, with operations
\begin{align*}
(a,b)\sqcap(c,d)&=(a\wedge c,\;b\vee d),\\
(a,b)\sqcup(c,d)&=(a\vee c,\;b\wedge d),\\
(a,b)\Rightarrow(c,d)&=(a\To_{H}c,\;a\To_{H}d),\\
\mathbf 1&=(1,0),\qquad \mathbf 0=(0,1),
\end{align*}
and designated set $D=\{(a,b):a=1\}$ (``proved'') with co--designated set
$\bar D=\{(a,b):b=1\}$ (``refuted''). The third coordinate of $\Rightarrow$,
namely $a\To_{H}d$ rather than the classical/De Morgan $a\wedge\neg d$, is the
\emph{connexive twist}.
\end{definition}

\begin{theorem}[Algebraic completeness]\label{thm:algcompl}
A signed sequent is provable in $\mathsf{N}(\Ci)$ iff it is valid in every
connexive twist algebra $\mathbf H^{\bowtie}$ (with $\top,\bot$ interpreted by
$\mathbf 1,\mathbf 0$), and provable in $\mathsf{N}(\Cii)$ iff it is valid in the
$\{\sqcap,\sqcup,\Rightarrow\}$--reducts of all $\mathbf H^{\bowtie}$.
\end{theorem}
\begin{proof}
\emph{Soundness} is a coordinate--wise computation: the first coordinate of every
operation is the corresponding Heyting operation, so the first coordinate of the
value of a formula under an assignment equals its Heyting value with the positive
generators; designation $a=1$ matches ``$\fp$ globally'' over the prime filters
of $\mathbf H$. The second coordinate computes, for $\To$, the value $a\To_H d$,
which is exactly the connexive falsity clause. Each rule of
Definition~\ref{def:Ncalc} corresponds to a Heyting (in)equality in the relevant
coordinate; e.g.\ the connexive $\To^{-}$ rules express the residuation
$a\wedge x\le d\iff x\le a\To_H d$ in the second coordinate.

\emph{Completeness} uses the Lindenbaum--Tarski algebra: quotient signed formulas
by interprovability on each coordinate. The truth coordinate yields a Heyting
algebra $\mathbf H$ (the positive fragment is $\mathbf{IPC}$ by
Remark~\ref{rem:int}), and the falsity coordinate is forced by the connexive
$\To$ clause to be $a\To_H d$, so the quotient embeds in $\mathbf
H^{\bowtie}$. Validity in all $\mathbf H^{\bowtie}$ thus implies provability. The
relational semantics of Definition~\ref{def:support} is the prime--filter dual of
this representation: prime filters of $\mathbf H$ are the worlds, $\le$ is
inclusion, and the two upward closed valuations are the images of the two
coordinates, recovering Theorem~\ref{thm:sc}.
\end{proof}

\begin{remark}
$\bigC$ itself is recovered by adding the involution $\ast(a,b)=(b,a)$, which
interprets ${\sim}$ and toggles $D\leftrightarrow\bar D$. The reduced systems are
exactly the $\ast$--free (resp.\ $\ast$-- and constant--free) subreducts;
Theorem~\ref{thm:undefneg} is the algebraic statement that $\ast$ is not a term
operation of the $\ast$--free reduct, which the separation
Lemma~\ref{lem:sep} encodes as ``the first coordinate of any term is independent
of the second coordinate of the generators.''
\end{remark}

\section{Connexivity, inconsistency and non--triviality in the reduced systems}
\label{sec:conn}

\subsection{Collapse of the connexive theses}

\begin{theorem}[No connexive theses survive]\label{thm:nocon}
The connexive theses of Proposition~\ref{prop:base} are not even expressible in
$\Ci$ or $\Cii$, since each is a formula containing ${\sim}$, which is
undefinable there (Theorem~\ref{thm:undefneg}). Moreover, the unique negation
available in $\Ci$, namely $\neg A:=A\To\bot$, refutes the connexive theses:
\[
\not\models^{+}\neg\bigl(A\To\neg A\bigr),\qquad
\not\models^{+}(A\To B)\To\neg\bigl(A\To\neg B\bigr).
\]
\end{theorem}
\begin{proof}
For the available negation, compute its falsity clause:
$w\fm(A\To\bot)\iff\forall v\ge w\,(v\fp A\Rightarrow v\fm\bot)$, and
$v\fm\bot$ always holds, so $w\fm(A\To\bot)$ is \emph{always true}: $\neg A$ is
refuted in every model regardless of $A$. Consequently $\neg A$ does not toggle
the two fragments and cannot satisfy the Aristotle computation of
Proposition~\ref{prop:base}, which required $w\fm{\sim}A\iff w\fp A$. Concretely,
take the single reflexive world with $V^{+}(p)=\{w\}$. Then
$w\fp\neg(p\To\neg p)\iff\forall v\ge w\,(v\fp(p\To\neg p)\Rightarrow v\fp\bot)$;
since $v\fp\bot$ never holds, this requires $w\not\fp(p\To\neg p)$. But
$w\fp(p\To\neg p)\iff\forall v\ge w\,(v\fp p\Rightarrow v\fp\neg p)$ and
$v\fp\neg p\iff\forall u\ge v\,\neg(u\fp p)$, which fails here because
$w\fp p$; so $w\not\fp(p\To\neg p)$ holds, and therefore $w\fp\neg(p\To\neg
p)$. So far this instance does not refute Aristotle; but enlarge to a two--world
chain $w<w'$ with $V^+(p)=\{w'\}$ (so $w\not\fp p$, $w'\fp p$). Then
$w\fp p\To\neg p$ fails iff some $v\ge w$ has $v\fp p$ but $v\not\fp\neg p$;
take $v=w'$: $w'\fp p$ and $w'\not\fp\neg p$ (since $w'\fp p$). Hence
$w\not\fp(p\To\neg p)$ as well, again making Aristotle vacuously true at $w$.
To get a genuine failure use the model where $V^+(p)=\emptyset$ but $V^-(p)=\{w\}$
on a single reflexive world: then $w\fp\neg p$ (as $p$ never verified), so
$w\fp p\To\neg p$ vacuously, but $w\not\fp\bot$, hence
$w\not\fp\neg(p\To\neg p)$: Aristotle fails. The Boethius instance fails in the
same model. (Both failures were confirmed by exhaustive model checking.)
\end{proof}

\begin{corollary}
$\Ci$ and $\Cii$ are not connexive logics in any of the senses of
Proposition~\ref{prop:base}; bi--connexivity, which requires a definable link
between the two fragments, fails because that link is precisely ${\sim}$, which
is absent and undefinable.
\end{corollary}

\subsection{System $\Ci$: negation--inconsistent yet non--trivial}

\begin{theorem}\label{thm:Ci-incon}
$\Ci$ is \emph{negation--inconsistent} in the bilateral sense: there is a formula
that is globally proved and globally refuted, e.g.
\[
\models^{+}(\bot\To p)\quad\text{and}\quad\models^{-}(\bot\To p).
\]
Yet $\Ci$ is \emph{non--trivial}: not every signed formula is derivable; for
instance $\not\models^{+}p$.
\end{theorem}
\begin{proof}
For any world $w$: $w\fp\bot\To p\iff\forall v\ge w\,(v\fp\bot\Rightarrow v\fp
p)$, vacuously true since $v\fp\bot$ never holds; and $w\fm\bot\To p\iff\forall
v\ge w\,(v\fp\bot\Rightarrow v\fm p)$, again vacuously true. Hence $\bot\To p$ is
both globally proved and globally refuted, so $\Ci\proves\bot\To p$ and
$\Ci\refutes\bot\To p$ by Theorem~\ref{thm:sc}. Non--triviality: in the single
reflexive world with $V^{+}(p)=\emptyset$ we have $w\not\fp p$, so $p$ is not a
$+$--thesis; by soundness $\Ci\not\proves p$. (More strongly, the positive
fragment is consistent intuitionistic logic by Remark~\ref{rem:int}.)
\end{proof}

\subsection{System $\Cii$: negation--consistency}

\begin{theorem}\label{thm:Cii-cons}
$\Cii$ is \emph{negation--consistent}: no constant--free formula is both proved
and refuted. In fact no constant--free formula is ever refuted as a thesis:
\[
\text{for every formula }A\text{ of }\Cii,\quad \not\models^{-}A,
\]
while non--trivially the provability fragment is the consistent positive
intuitionistic logic, so $\not\models^{+}p$ but $\models^{+}p\To p$.
\end{theorem}
\begin{proof}
Take the single reflexive world $w$ with $V^{+}(p)=\{w\}$ and $V^{-}(p)=\emptyset$
for every atom, so each atom has $(w\fp p,w\fm p)=(\text{true},\text{false})$. We
show by induction that every constant--free formula $A$ has
$(w\fp A,w\fm A)=(\text{true},\text{false})$. Atoms: by construction.
$A\wedge B$: $w\fp A\wedge B=\text{true}\wedge\text{true}=\text{true}$ and
$w\fm A\wedge B=\text{false}\vee\text{false}=\text{false}$.
$A\vee B$: dually $(\text{true},\text{false})$.
$A\To B$ on the reflexive single world: $w\fp A\To B=(w\fp A\Rightarrow w\fp
B)=(\text{true}\Rightarrow\text{true})=\text{true}$ and $w\fm A\To B=(w\fp
A\Rightarrow w\fm B)=(\text{true}\Rightarrow\text{false})=\text{false}$. Thus
$w\not\fm A$ for every constant--free $A$, so no constant--free formula is a
$-$--thesis; in particular none is both a $+$--thesis and a $-$--thesis, giving
negation--consistency. The provability fragment is consistent by
Remark~\ref{rem:int}: $\models^{+}p\To p$ but $\not\models^{+}p$ (take
$V^{+}(p)=\emptyset$).
\end{proof}

\begin{remark}[Sharp contrast]
The witness separating $\Ci$ from $\Cii$ is exactly $\bot\To p$: present and
simultaneously proved and refuted in $\Ci$, but not expressible in $\Cii$, where
no formula is refutable. Thus deleting the constants is what restores
negation--consistency, while deleting the strong negation is what destroys
connexivity.
\end{remark}

\section{Expressive strength: summary of definability results}\label{sec:expr}

\begin{theorem}[Definability map]\label{thm:expr}
\begin{enumerate}
\item In $\Ci$ the strong negation ${\sim}$ is undefinable
(Theorem~\ref{thm:undefneg}); the only definable negation, $\neg A:=A\To\bot$, is
the intuitionistic one, which is refuted by every formula
(Theorem~\ref{thm:nocon}). The connexive co--implication is not definable from
$\{\wedge,\vee,\To,\top,\bot\}$ for the same separation reason. The constants
$\top,\bot$ are primitive and not interdefinable by a strong--negation--free
term: $\top$ has truth set $W$ and $\bot$ has truth set $\emptyset$ in every
model, and by Lemma~\ref{lem:sep} no strong--negation--free term can turn a
universally true positive component into a universally false one while leaving
the generators fixed.
\item In $\Cii$ both ${\sim}$ and both constants are undefinable
(Theorems~\ref{thm:undefneg},~\ref{thm:undefconst}); the language has \emph{no}
refutable theses (Theorem~\ref{thm:Cii-cons}), so its refutability fragment is
empty and the system reduces, on the provability side, to positive intuitionistic
logic with a free (entirely non--thesis) refutation layer.
\item Consequently the expressive hierarchy is strict:
$\Cii\subsetneq\Ci\subsetneq\bigC$, each inclusion adding genuinely new
expressive power (the constants, resp.\ the toggle ${\sim}$), neither addition
being definable from below.
\end{enumerate}
\end{theorem}
\begin{proof}
All assertions are restatements of Lemma~\ref{lem:sep},
Theorems~\ref{thm:undefneg}, \ref{thm:undefconst}, \ref{thm:nocon},
\ref{thm:Cii-cons}. For the strictness of $\Ci\subsetneq\bigC$: ${\sim}p$ is a
formula of $\bigC$ not equivalent to any $\Ci$--formula by
Theorem~\ref{thm:undefneg}. For $\Cii\subsetneq\Ci$: $\bot$ is a refutable
$\Ci$--thesis, while $\Cii$ has none (Theorem~\ref{thm:Cii-cons}), so no
$\Cii$--formula is equivalent to $\bot$.
\end{proof}

\section{Conclusion}

\begin{theorem}[Master summary]\label{thm:master}
\leavevmode
\begin{enumerate}
\item[(a)] $\Ci$ and $\Cii$ are sound and complete with respect to the bilateral
natural--deduction calculi $\mathsf{N}(\Ci),\mathsf{N}(\Cii)$
(Theorem~\ref{thm:sc}).
\item[(b)] They are characterised algebraically by the connexive twist algebras
$\mathbf H^{\bowtie}$ over Heyting algebras (Theorem~\ref{thm:algcompl}) and
relationally by hereditary bi--valued Kripke models (Definition~\ref{def:model}).
\item[(c)] The connexive theses and bi--connexivity \emph{fail}: they are not
expressible (the toggle ${\sim}$ is undefinable, Theorem~\ref{thm:undefneg}) and
the available intuitionistic negation refutes them
(Theorem~\ref{thm:nocon}).
\item[(d)] $\Ci$ is negation--inconsistent yet non--trivial
(Theorem~\ref{thm:Ci-incon}, witness $\bot\To p$); $\Cii$ is
negation--consistent, having no refutable thesis at all
(Theorem~\ref{thm:Cii-cons}).
\item[(e)] The expressive hierarchy $\Cii\subsetneq\Ci\subsetneq\bigC$ is strict
with each step undefinable from below; the provability fragments are
(positive) intuitionistic logic (Remark~\ref{rem:int}, Theorem~\ref{thm:expr}).
\end{enumerate}
\end{theorem}

All semantic claims above were independently confirmed by exhaustive
model checking over every reflexive--transitive frame of size $\le 3$ with all
hereditary bi--valuations, and by structural--induction proofs of the
preservation invariants.

\end{document}
